% !TeX root = ../main.tex
% Chapter 1 — Introduction
% Self-contained section file: no preamble, no \documentclass.
% Assumes the main document loads: graphicx, booktabs, tikz (libraries: positioning, arrows.meta, fit, backgrounds), hyperref, and a citation backend (biblatex/natbib) resolving the keys below.

\chapter{Introduction}
\label{ch:introduction}

\section{Background and significance}
\label{sec:intro-background}

Cats are obligate maskers of pain. As a prey-adapted species they suppress overt
distress behaviours in the presence of an observer, so that the very act of
clinical examination perturbs the signal a clinician is trying to read. This
makes feline acute-pain assessment one of the harder problems in companion-animal
welfare: the patient cannot self-report, the behavioural channel is actively
concealed, and the physiological channels (heart rate, blood pressure) are
confounded by the stress of handling. The consequence is systematic
under-treatment of pain in cats relative to dogs, a gap that the veterinary
literature has repeatedly identified as a welfare and analgesia-stewardship
concern.

The dominant response to this problem has been to ground pain assessment in
\emph{facial} rather than \emph{behavioural} expression. Facial movement is far
harder to suppress voluntarily than posture or vocalisation, and it can be coded
objectively through anatomically-defined muscle movements. The Feline Facial
Action Coding System (CatFACS), an adaptation of the human Facial Action Coding
System (FACS), decomposes the cat face into discrete \emph{action units} (AUs)
that correspond to identifiable muscle actions. Building on this foundation, the
\emph{Feline Grimace Scale} (FGS) operationalises pain into five ordinal action
units---ear position, orbital tightening, muzzle tension, whisker change, and
head position---each scored on a three-level ordinal scale
($0 = \text{absent}$, $1 = \text{moderate or uncertain}$, $2 = \text{obvious}$).
The five per-AU scores are summed to a $0$--$10$ composite, and a normalised
ratio threshold of $>0.39$ separates the painful from the non-painful state
\autocite{evangelista2019}. The FGS is attractive precisely because it is cheap,
fast, and rater-trainable; its original validation reported strong discrimination
of acute pain (area under the curve $\approx 0.94$, sensitivity $\approx 90.7\%$,
specificity $\approx 86.6\%$ at the $>0.39$ cut-point) against a clinical
reference \autocite{evangelista2019}. Figure~\ref{fig:fgs-five-au} summarises the
five-AU structure that the present work treats as its target ontology.

% ---------------------------------------------------------------------------
% Figure 1.1 — FGS five-AU diagram (captioned placeholder, TikZ schematic).
% ---------------------------------------------------------------------------
\begin{figure}[t]
  \centering
  \begin{tikzpicture}[
      font=\small,
      au/.style={draw, rounded corners, minimum width=2.5cm, minimum height=1.0cm,
                 align=center, fill=gray!8},
      lvl/.style={font=\scriptsize\itshape, gray},
      every node/.style={align=center},
    ]
    % central face placeholder
    \node[draw, circle, minimum size=2.2cm, fill=gray!4] (face) {cat face\\(crop)};
    % five AUs around the face
    \node[au, above left=0.9cm and 1.6cm of face]  (ear)      {Ear position};
    \node[au, above right=0.9cm and 1.6cm of face] (orbital)  {Orbital\\tightening};
    \node[au, right=2.0cm of face]                 (muzzle)   {Muzzle\\tension};
    \node[au, below right=0.9cm and 1.6cm of face] (whisker)  {Whisker\\change};
    \node[au, below left=0.9cm and 1.6cm of face]  (head)     {Head\\position};
    % connectors
    \foreach \n in {ear,orbital,muzzle,whisker,head}{
      \draw[-{Latex[length=2mm]}, gray] (face) -- (\n);
    }
    % per-AU ordinal scale annotation
    \node[lvl, below=0.05cm of ear]     {$0\,/\,1\,/\,2$};
    \node[lvl, below=0.05cm of orbital] {$0\,/\,1\,/\,2$};
    \node[lvl, below=0.05cm of muzzle]  {$0\,/\,1\,/\,2$};
    \node[lvl, below=0.05cm of whisker] {$0\,/\,1\,/\,2$};
    \node[lvl, below=0.05cm of head]    {$0\,/\,1\,/\,2$};
    % aggregation note
    \node[draw, dashed, rounded corners, below=1.4cm of face, fill=gray!4,
          minimum width=8.2cm, align=center] (sum)
      {five per-AU scores $\;\rightarrow\;$ sum $\in\{0,\dots,10\}$
       $\;\rightarrow\;$ normalised ratio $>0.39$ flags \emph{pain}};
    \draw[-{Latex[length=2mm]}, gray] (head) |- (sum);
  \end{tikzpicture}
  \caption[The five FGS action units]{Schematic of the Feline Grimace Scale
    \autocite{evangelista2019}. Each of the five action units is scored on a
    three-level ordinal scale ($0/1/2$); the five scores are summed to a $0$--$10$
    composite and converted to a normalised ratio, with $>0.39$ flagging pain.
    This diagram fixes the target ontology used throughout the thesis; it is a
    structural schematic, not an example annotation. (Placeholder figure.)}
  \label{fig:fgs-five-au}
\end{figure}

A recent wave of machine-learning work has sought to automate this pipeline.
Two methodological families dominate. The first is \emph{landmark-then-geometry}:
a facial-landmark detector locates anatomical points, hand-engineered geometric
features are derived, and a classifier (frequently a gradient-boosted tree)
predicts pain or grimace score. Representative of this family,
\textcite{steagall2023} report an automated FGS pipeline reaching $95.5\%$
agreement using a minimum-aggregation rule over the whisker and head action units,
and \textcite{feighelstein2023} report a binary cat-pain classifier in which a
landmark-based model ($\approx 77\%$) outperformed a direct deep-learning
classifier ($\approx 65\%$), alongside an explainability analysis ranking the
mouth region above the eyes and ears in feature importance. The second family
relies on facial-landmark corpora to standardise the geometry: CatFLW provides
$48$ CatFACS-aligned landmarks per face with reported normalised mean error in
the $9$--$26\%$ range at roughly seven landmark-points of cost
\autocite{martvel2024catflw}. In parallel, foundation vision encoders such as
DINOv2 \autocite{oquab2023dinov2} have made it cheap to extract strong,
general-purpose facial features without task-specific pre-training, lowering the
data and compute floor for small-corpus ordinal modelling. Most recently, a COSMIN
scoping review of feline pain instruments has begun to systematise the
measurement-quality evidence for \emph{human-rater} instruments
\autocite{leesteagall2026}; automated scoring falls outside that review's topic
scope, so the question of what \emph{measurement} guarantees an automated scorer
should carry remains open.

\section{The gap, in our terms}
\label{sec:intro-gap}

The automated-FGS and binary-cat-pain literature, summarised above, is impressive
on its own terms---but those terms are accuracy on a single-source corpus. We
identify three gaps that this prior work does not address, and which together
motivate the present thesis.

\paragraph{(G1) No transportable trustworthiness method.}
Prior automated work optimises a point accuracy (or AUC) on one corpus and reports
it. What it does not produce is a \emph{method}---a piece of released, runnable
machinery that a future investigator can carry to the next corpus and re-run to
obtain a trustworthiness measurement. The headline number in
\textcite{steagall2023} or \textcite{feighelstein2023} is a fact about
\emph{their} dataset; it is not an instrument that travels. Because no open,
graded cat-FGS dataset exists, a contribution that is merely ``high accuracy on
corpus $X$'' is, for a feline-pain scorer, dead on arrival under independent
replication: a reviewer can dismiss it as a backbone swap on inaccessible data.
The portable artefact---the reusable measurement protocol---is the missing
deliverable.

\paragraph{(G2) No audit of capture-condition confounds.}
A facial-pain corpus is acquired under conditions (lighting, camera, restraint,
clinical context) that may correlate with the pain label itself. If brightness,
blur, aspect ratio, or background can predict the pain label, then a classifier
can score well by reading acquisition context rather than facial expression, and
the reported accuracy is an artefact of capture rather than evidence of pain
detection. Prior work does not systematically audit this entanglement, and---more
importantly---does not ship the audit as a protocol that the next corpus can run.

\paragraph{(G3) Low-prevalence measurement-error propagation is ignored.}
Acute clinical pain is a low-prevalence event in many realistic deployment
settings. At low prevalence, even small per-sample labelling and classification
errors propagate non-linearly into the operating characteristics that a clinician
actually relies upon. \textcite{chavoshi2025} make this concrete: at a $10\%$
prevalence, an instrument whose raw error looks benign can collapse to roughly
$53\%$ effective sensitivity once measurement error is propagated through the base
rate. A scorer reported only as a single accuracy figure, with no confidence
interval and no prevalence-aware operating point, silently hides this failure
mode.

These three gaps share a common root: the field reports \emph{performance} but not
\emph{trustworthiness}, and it reports both as corpus-bound numbers rather than
transportable methods. The present thesis is organised around closing that gap
without pretending to data we do not have.

\section{Objectives}
\label{sec:intro-objectives}

The thesis pursues four numbered objectives, ordered from the headline
\emph{portable method} down to the conceded engineering substrate.

\begin{enumerate}
  \item \textbf{Build a portable, power-aware confound-attribution protocol.}
    Construct a reusable, one-directional capture-condition audit (the
    FGS--BG--Gap, the per-AU energy-based pointing-game read as
    saliency-as-confound-evidence, EBPG, and a VLM judge-bias probe) that any
    future facial-pain-scorer corpus can run, at a pre-registered power, to
    detect---never to ``rule out''---entanglement between the pain label and
    acquisition context. Reporting is one-directional at this
    power~\autocite{adebayo2022posthoc}: the only honest conclusion is ``no
    confound detected at this power.'' This is the headline contribution, and it
    is validated today on planted positive and negative controls.
  \item \textbf{Provide a guarded VLM-as-AU-rater reliability check.} As a
    subordinate, inspected-not-validated reliability check, determine whether
    a frozen vision--language model (VLM) can weak-label the five FGS action units
    at \emph{human-rater agreement}---a $\kappa$ that measures weak-labelling
    capability only when the vet anchor is independent and the rubric handed to
    the VLM differs from the rubric the vet scored from, and otherwise measures
    rubric-following, so each run must state which it measures---quantified as a per-AU quadratic-weighted
    Cohen's $\kappa$ against a veterinary anchor, reported with a confidence-interval
    lower bound. The CI-lower-bound gate is textbook clinimetrics and the
    rubric-paraphrase guard is published practice; neither is claimed as novel.
    This check is retained as kill-tree insurance---the sole-survivor headline
    should the confound leg degrade---not as a co-equal second pillar.
  \item \textbf{Ship a welfare-asymmetric binary-plus-abstention decision floor.}
    Deliver a calibrated binary pain decision at a \emph{fixed} high-sensitivity
    operating point (pain-recall $\geq 0.90$, anchored to
    \textcite{evangelista2019}; never a symmetric Youden-$J$/$F_1$ knee), made
    visible through a welfare-asymmetric decision curve, with a defer-to-vet
    abstention boundary reported as a one-sided $95\%$ negative-predictive-value
    lower bound.
  \item \textbf{Build---but only inspect, not validate---a $0$--$10$ ordinal
    engine.} Implement a frozen DINOv2 encoder feeding five per-AU ordinal heads
    and a distributional decode to a $0$--$10$ composite, and ship it as an
    \emph{inspected-not-validated} exploratory layer. This engine is conceded as
    plumbing; it is never advanced as a validated graded claim.
\end{enumerate}

Objective~1 is the headline transportable contribution; Objective~2 is the
guarded reliability check that rides below it and serves as kill-tree insurance;
Objective~3 is the cited trustworthiness wrapper, an honest decision floor that
stands on its own even if the graded layer is dropped, but is supporting
evidence rather than a headline contribution; Objective~4 is built for
completeness and future extension but is firewalled out of every validated
results table. The contribution map is summarised in
Figure~\ref{fig:contribution-map}.

% ---------------------------------------------------------------------------
% Figure 1.2 — Thesis contribution map (TikZ).
% ---------------------------------------------------------------------------
\begin{figure}[t]
  \centering
  \begin{tikzpicture}[
      font=\small,
      box/.style={draw, rounded corners, minimum height=1.0cm, align=center,
                  text width=3.7cm, inner sep=4pt},
      head/.style={box, fill=gray!14, very thick},
      guard/.style={box, fill=gray!8},
      frame/.style={box, fill=gray!6},
      engine/.style={box, dashed, fill=white},
      lbl/.style={font=\scriptsize\itshape, gray},
    ]
    % headline portable method (left) + guarded subordinate check (right)
    \node[head] (m1) {\textbf{Headline method}\\Power-aware confound-\\attribution protocol\\(FGS-BG-Gap + per-AU\\EBPG + VLM judge-bias)};
    \node[guard, right=1.0cm of m1] (m2) {\emph{Guarded check}\\VLM-as-AU-rater\\per-AU quadratic $\kappa$\\vs.\ vet (CI-LB-gated;\\inspected-not-validated)};
    \node[lbl, above=0.15cm of m1.north east] {headline portable method (left) + guarded subordinate $\kappa$ check (right)};
    % the wrapper frame underneath
    \node[frame, below=1.2cm of m1, text width=3.7cm] (op)
      {Welfare-asymmetric decision curve\\@ pain-recall $\geq 0.90$};
    \node[frame, below=1.2cm of m2, text width=3.7cm] (ab)
      {One-sided $95\%$ NPV\\lower-bound abstention};
    \node[lbl, below=0.1cm of op.south east] {trustworthiness wrapper --- cited supporting plumbing};
    % conceded engine at the base, centred under the two wrapper frames
    \node[engine, below=1.5cm of op.south east, anchor=north, text width=8.4cm] (eng)
      {\textbf{Conceded engine (plumbing --- never claimed novel):}
       frozen DINOv2 ViT-S/14 (register variant \texttt{dinov2\_vits14\_reg} the
       field default) $\rightarrow$ 5 per-AU CORN heads $\rightarrow$
       $0/1/2 \rightarrow$ sum $0$--$10 \rightarrow$ distributional decode.
       Graded layer ships \emph{inspected-not-validated}.};
    % connectors
    \draw[-{Latex[length=2mm]}, gray] (eng) -- (op);
    \draw[-{Latex[length=2mm]}, gray] (eng) -- (ab);
    \draw[-{Latex[length=2mm]}, gray] (op) -- (m1);
    \draw[-{Latex[length=2mm]}, gray] (ab) -- (m2);
  \end{tikzpicture}
  \caption[Thesis contribution map]{Contribution map. The single headline
    deliverable is a portable, power-aware capture-condition
    confound-attribution protocol (top left). The VLM-as-AU-rater per-AU
    $\kappa$ protocol (top right) is a guarded, inspected-not-validated
    reliability check that rides below it as kill-tree insurance, not a co-equal
    second pillar. The trustworthiness wrapper (middle) supplies the
    welfare-asymmetric operating point and the abstention floor and is cited
    supporting plumbing rather than a headline. The DINOv2$+$CORN engine
    (bottom) is conceded plumbing and is never claimed as a contribution; its
    graded output ships inspected-not-validated. (Placeholder schematic.)}
  \label{fig:contribution-map}
\end{figure}

\section{Scope and honesty boundary}
\label{sec:intro-scope}

This thesis is, at the time of writing, a \emph{research proposal and methodology}
backed by a build specification; it reports an evaluation protocol, gate-conditioned
expected outcomes, and decision rules, rather than completed experimental results.
No accuracy, $\kappa$, or calibration number reported as ``ours'' appears in this
document, because no such experiment has yet been run. Literature numbers
(e.g.\ those of \textcite{evangelista2019}, \textcite{steagall2023},
\textcite{feighelstein2023}, \textcite{martvel2024catflw}, and
\textcite{chavoshi2025}) are reproduced only as background and are explicitly
attributed to their sources.

Several scope boundaries follow from honesty about what the project can and cannot
deliver, and are held as hard constraints throughout:

\begin{itemize}
  \item \textbf{Single corpus, single anchor.} The empirical substrate is a single
    forked Roboflow cat-pain corpus (binary pain / no\_pain boxes,
    $\sim 2040$ records / $\sim 1547$ distinct images, $\sim 12$--$13\%$
    prevalence, with a leaky shipped split that is discarded) plus one veterinary
    anchor whose size is fixed by an up-front power calculation. No open, graded
    cat-FGS dataset exists; CatFLW \autocite{martvel2024catflw} is used for
    eye-alignment landmarks only, not as a graded anchor, and an open
    horse-grimace set is used solely as a decode unit-test fixture (no horse
    number ever appears in a results table).
  \item \textbf{Solo compute.} The work is built by a single developer on an Apple
    M4 / MPS machine (no local CUDA); detector training alone is offloaded to a
    cloud T4, and the frozen encoder plus the ordinal heads run locally on cached
    features.
  \item \textbf{Decision-support framing only.} The output is a triage
    recommendation---``grimace consistent with pain, $X/10$; recommend veterinary
    assessment''---and \emph{never} an autonomous analgesia trigger. The negative
    class is documented as an unknown mixture (possibly sedated or post-operative
    animals), not as a clean ``no-pain'' reference.
  \item \textbf{Engine conceded.} The DINOv2$+$CORN engine is treated as plumbing
    throughout and is never claimed as a novel contribution; the word
    ``graded'' is struck from every validated-claim sentence, and the thesis
    abstract holds with the graded output dropped.
  \item \textbf{Anti-benchmark discipline.} The thesis never claims to ``beat''
    the $77\%$ / $79\%$ / $95\%$ figures of prior work. Every reported rate is
    accompanied by its distinct-pain-cat denominator and a Clopper--Pearson or
    bootstrap confidence interval, and accuracy is treated as an anti-benchmark
    rather than the objective.
\end{itemize}

\section{Research procedure}
\label{sec:intro-procedure}

The work is executed as a strict, blocking \emph{gate order}: each gate writes an
auditable artefact and blocks all downstream quantitative work until it passes.
The order is risk-first---the cheapest checks that can kill or reframe the project
run before any veterinary hour or GPU spend. The eight gates are summarised below
and depicted in Figure~\ref{fig:gate-flow}.

\begin{enumerate}
  \item \textbf{G0 --- Power calculations and vet-budget pre-registration.}
    Fix the single integer (how many faces the vet scores) by three zero-data
    power calculations: the faces needed for a per-AU $\kappa$ half-width
    $\leq 0.15$; the faces needed to certify a useful negative-predictive-value
    band at an acceptable abstention rate; and whether the $>0.39$ operating-point
    interval is reportable at the budgeted $n$. Blocks everything quantitative.
  \item \textbf{G1 --- Per-cat merge.} Collapse clips to true individuals,
    validated against trusted camera ids, so that no individual straddles a fold
    and the honest distinct-pain-cat denominator can be printed.
  \item \textbf{G2 --- Capture-condition confound audit.} Run the one-directional
    portable audit (Objective~2) before any spend; it can kill or reframe the
    corpus, and it publishes either way.
  \item \textbf{G3 --- Frozen, hashed, cat-disjoint hold-out.} Freeze and hash the
    test manifest so that a continuous-integration check aborts any run that can
    read it---structurally defeating post-hoc goalpost moves.
  \item \textbf{G4 --- MPS compute-correctness pre-checks (blocking).} Assert
    MPS-versus-CPU encoder logit parity and pass a synthetic ordinal
    decode-to-sum-to-threshold unit test, so that no number from silently-wrong
    hardware logits is ever believed.
  \item \textbf{G5 --- Alignment / normalised-mean-error gate.} Check that the
    detector crop is geometrically clean enough to feed the ordinal heads, adding
    a two-point eye-similarity alignment if not, so that ``calibration'' never
    silently measures crop quality.
  \item \textbf{G1-B --- Weak-label reliability $\kappa$ pilot.} Score the per-AU
    VLM-versus-vet quadratic $\kappa$ on the anchor (Objective~1); the pilot itself
    selects the VLM, and the gate fires on the $\kappa$ confidence-interval
    \emph{lower bound}, not the point estimate.
  \item \textbf{G6 --- Severity-cell count gate.} If the anchor yields only
    single-digit AU$=2$ (severe) cells, collapse the high end of the scale by a
    pre-committed rule rather than printing an uninformative per-cell number.
\end{enumerate}

% ---------------------------------------------------------------------------
% Figure 1.3 — Gate flow (TikZ).
% ---------------------------------------------------------------------------
\begin{figure}[t]
  \centering
  \begin{tikzpicture}[
      font=\scriptsize,
      gate/.style={draw, rounded corners, minimum width=1.65cm, minimum height=0.85cm,
                   align=center, fill=gray!8},
      block/.style={gate, fill=gray!18, very thick},
      flow/.style={-{Latex[length=1.6mm]}, gray},
    ]
    \node[gate] (g0) {G0\\power\\+ budget};
    \node[gate, right=0.55cm of g0] (g1) {G1\\per-cat\\merge};
    \node[gate, right=0.55cm of g1] (g2) {G2\\confound\\audit};
    \node[gate, right=0.55cm of g2] (g3) {G3\\frozen\\hold-out};
    \node[block, below=0.7cm of g3] (g4) {G4\\MPS\\(blocking)};
    \node[gate, left=0.55cm of g4] (g5) {G5\\align\\/ NME};
    \node[gate, left=0.55cm of g5] (g1b) {G1-B\\$\kappa$ pilot\\(CI-LB)};
    \node[gate, left=0.55cm of g1b] (g6) {G6\\severity\\collapse};
    \draw[flow] (g0) -- (g1);
    \draw[flow] (g1) -- (g2);
    \draw[flow] (g2) -- (g3);
    \draw[flow] (g3) -- (g4);
    \draw[flow] (g4) -- (g5);
    \draw[flow] (g5) -- (g1b);
    \draw[flow] (g1b) -- (g6);
  \end{tikzpicture}
  \caption[Gate run-order]{The blocking gate run-order
    (G0 $\rightarrow$ G1 $\rightarrow$ G2 $\rightarrow$ G3 $\rightarrow$ G4
    $\rightarrow$ G5 $\rightarrow$ G1-B $\rightarrow$ G6). Each gate emits an
    auditable artefact and blocks downstream quantitative work; G4 (MPS
    compute-correctness) is hard-blocking. The order is risk-first: the cheapest
    kill-or-reframe checks run before any veterinary or GPU spend.}
  \label{fig:gate-flow}
\end{figure}

A kill/pivot rule is attached to the gate order: if the orbital, ear, or head
$\kappa$ lower bound at G1-B falls below its pre-registered floor, the graded
layer is dropped entirely and the project ships the calibrated
binary-plus-abstention floor with the confound protocol; the headline confound
protocol still stands, because it is validated independently of the $\kappa$
result. Conversely, the $\kappa$ check is retained as kill-tree insurance: should
the confound leg itself degrade, the $\kappa$ check---once the independent per-AU
vet anchor exists---becomes the sole-survivor headline, because it measures a
real fact whether the result is strong or mediocre.

\section{Expected benefits and contributions}
\label{sec:intro-contributions}

The thesis is designed so that its value does not depend on any single experiment
firing favourably. Its expected contributions are one headline released,
dataset-agnostic method, a guarded reliability check below it, and a cited
supporting wrapper:

\begin{enumerate}
  \item \textbf{A portable, power-aware capture-condition confound-attribution
    protocol (the headline).} The one-directional FGS--BG--Gap plus per-AU
    EBPG-as-confound-evidence audit, augmented by a VLM judge-bias probe and run
    at a pre-registered power, runnable on any future facial-pain-scorer corpus,
    whose deliverable is the protocol itself rather than a verdict about one
    dataset. It is validated today on planted positive and negative controls and
    reports one-directionally at this power~\autocite{adebayo2022posthoc}. The
    claimed novelty is the assembly, the per-AU
    EBPG-as-confound-evidence reading, and the instantiation of
    equivalence-style audit hygiene---never any individual primitive or the
    underlying statistics.
  \item \textbf{A guarded VLM-as-AU-rater $\kappa$ reliability check (subordinate;
    result pending the independent vet anchor).} Released code that
    scores any face corpus's per-AU VLM labels against a veterinary anchor and
    reports five quadratic-weighted $\kappa$ values with confidence-interval lower
    bounds. The forked corpus's numbers merely \emph{instantiate} the check;
    the deliverable is the reusable machinery; the number it will
    eventually produce---whether a frozen VLM weak-labels feline FGS action
    units at human-rater agreement---remains pending the independent per-AU
    vet anchor (Gate 0), since the current binary corpus cannot yield $0/1/2$
    AU ground truth. The CI-lower-bound gate is textbook clinimetrics and the
    rubric-paraphrase guard is published; neither is a novel increment. This check
    is kill-tree insurance, not a co-equal headline.
  \item \textbf{A welfare-asymmetric decision-curve and abstention floor (cited
    supporting plumbing).} An
    operating point fixed at high sensitivity with the undertreat-to-overtreat
    harm ratio swept as a range, plus a defer-to-vet boundary reported with
    finite-sample one-sided lower bounds. This is the clinical-decision
    component that distinguishes the trustworthiness wrapper from generic
    calibration, and it makes the low-prevalence error-propagation concern of
    \textcite{chavoshi2025} an explicit, reported quantity rather than a hidden
    failure mode. It is supporting evidence, not a headline contribution.
\end{enumerate}

Taken together these artefacts answer the three gaps of
Section~\ref{sec:intro-gap} without overclaiming: they supply the transportable
trustworthiness method (G1), the capture-condition audit (G2), and the
prevalence-aware operating point (G3) that the prior automated-FGS literature
leaves open. The remainder of the thesis develops each in turn---the FGS and
automated-scoring background and the measurement-quality literature
(Chapter~2), the methodology and gate-bound evaluation protocol, and the
gate-conditioned expected outcomes and decision rules---always under the standing
constraint that the DINOv2$+$CORN engine is conceded plumbing and that no graded
claim is ever validated.
